\documentclass{article}
\usepackage{microtype}
\usepackage{graphicx}
\usepackage{booktabs}
\usepackage{hyperref}
\usepackage[preprint]{icml2026}
\usepackage{amsmath,amssymb}
\usepackage[T1]{fontenc}
\icmltitlerunning{Selecting Whole Benchmarks to Predict Unmeasured Scores}
\begin{document}
\twocolumn[
\icmltitle{Selecting Whole Benchmarks to Predict Unmeasured Scores:\\Conditional Evidence from Public Leaderboards}
\begin{icmlauthorlist}
\icmlauthor{ResearchWorld Study Team}{rw}
\end{icmlauthorlist}
\icmlaffiliation{rw}{Independent computational research study}
\icmlcorrespondingauthor{ResearchWorld Study Team}{via the accompanying repository issue tracker}
\icmlkeywords{benchmark selection, score prediction, model evaluation, incomplete data}
\vskip 0.3in
]
\printAffiliationsAndNotice{}
\input{paper/body.tex}
\clearpage
\section*{Impact Statement}
This study aims to reduce the measurements needed for model evaluation using already public aggregate scores. Predicted benchmark scores should be labelled as predictions, with the measured query set and uncertainty conditions disclosed. Presenting them as actual evaluations could mislead model users, obscure unusual weaknesses, or reward models that resemble the historical leaderboard population. Historical coverage favors widely evaluated models and selected languages; performance prediction does not replace direct safety, fairness, or domain-specific evaluation. The study neither releases model weights nor collects new personal data. Its empirical savings concern benchmark counts; any environmental or economic savings from future deployment require separate cost measurements.
\bibliography{context/literature/references}
\bibliographystyle{icml2026}
\clearpage
\appendix
\onecolumn
\input{paper/appendix.tex}
\end{document}
